\documentclass[11pt,reqno]{amsart}
\usepackage[T1]{fontenc}
\usepackage{amsmath,amssymb,amsthm}
\usepackage[margin=1in]{geometry}
\usepackage[colorlinks=true,linkcolor=blue]{hyperref}
\newtheorem{theorem}{Theorem}[section]
\newtheorem{lemma}[theorem]{Lemma}
\newtheorem{proposition}[theorem]{Proposition}
\newtheorem{corollary}[theorem]{Corollary}
\theoremstyle{definition}
\newtheorem{definition}[theorem]{Definition}
\newtheorem{hypothesis}[theorem]{Hypothesis}
\theoremstyle{remark}
\newtheorem{remark}[theorem]{Remark}
\emergencystretch=3em
\title{REV-PLAN: exact shared interfaces for the AKLT gap revision}
\date{}
\begin{document}
\maketitle
\noindent These are mathematical statements for parallel proof writing.
No proof or certificate verification is asserted by this interface file.
All expectations and functionals below are finite sums.
The statements with hypotheses are internal implications; the target
main theorem is unconditional after the computational input theorems
are proved. Labels are shared with the complete inventory in PLAN.md.


\section{Result and conventions}

% Owner: A
\begin{definition}\label{d:model}
For integers \(L,N\ge3\), let \(\Lambda=\Lambda_{L,N}=\mathbb Z_L\times\mathbb Z_N\)
with its simple, unoriented nearest-neighbour edge set \(E_{L,N}\).
The Hilbert space is \(\mathcal H_\Lambda=\bigotimes_{x\in\Lambda}V_4\).
Each site carries the spin-\(2\) space \(V_4\), and \(P_{\{x,y\}}\)
is the orthogonal projection onto total spin \(4\) in \(V_4\otimes V_4\).
Fix one orientation of each edge and set
\[
 q_{xy}=z_{x,1}z_{y,2}-z_{x,2}z_{y,1},\qquad
 B_{L,N}=\prod_{e\in E_{L,N}}q_e,\qquad
 H_{L,N}=\sum_{e\in E_{L,N}}P_e.
\]
Every edge is counted once. The spectral gap of \(H_{L,N}\) means its
smallest positive eigenvalue.
\end{definition}


\section{Finite moment algebra and the response input}

% Owner: C
\begin{definition}\label{d:moment}
Write \(\mathbb S^2=\{a\in\mathbb R^3:|a|=1\}\), put
\(\phi=(1+\sqrt5)/2\), \(r=\sqrt{\phi+2}\), and set
\[
 D=\frac1r\{(0,\pm1,\pm\phi),(\pm1,\pm\phi,0),(\pm\phi,0,\pm1)\},
 \qquad
 \mathcal M(f)=\frac1{12}\sum_{a\in D}f(a).
\]
All indicated signs are independent. Define finite probability laws
\(\lambda_{\mathrm U}=\frac1{12}\sum_{a\in D}\delta_a\) and
\(\lambda_{\mathrm D}=\frac35\lambda_{\mathrm U}+\frac25\delta_0\).
Products of these functionals always mean finite sums.
For plane sites, \(d\) denotes ambient \(\ell^1\) distance and
\(\operatorname{dist}_\infty\) denotes ambient \(\ell^\infty\) distance.
Distances to an empty set are \(+\infty\); a product over an empty site
set has its one empty configuration. For finite probabilities, use
\[
 d_{\rm TV}(p,p')=\tfrac12\sum_z|p(z)-p'(z)|,\qquad
 \operatorname{Var}_\mu f=\mu(|f-\mu f|^2),\qquad
 d_H(n,n')=\sum_x\mathbf1_{\{n_x\ne n'_x\}} .
\]
\end{definition}

% Owner: C
\begin{lemma}\label{l:moments}
The set \(D\) has twelve distinct unit vectors and is antipodal.
For \(i,j,k,l\in\{1,2,3\}\),
\[
 \begin{gathered}
 \mathcal M(1)=1,\qquad \mathcal M(n_i)=\mathcal M(n_in_jn_k)=0,\qquad
 \mathcal M(n_in_j)=\delta_{ij}/3,\\
 \mathcal M(n_in_jn_kn_l)=
 (\delta_{ij}\delta_{kl}+\delta_{ik}\delta_{jl}+\delta_{il}\delta_{jk})/15.
 \end{gathered}
\]
For every polynomial \(p\) of total degree at most four and every
\(R\in\mathrm O(3)\), \(\mathcal M(p\circ R)=\mathcal M(p)\).
The same covariance holds for \(\lambda_{\mathrm D}\); its constant
moment is one and its moments of positive degree at most four are
\(3/5\) times those displayed above.
These assertions hold site by site for polynomials of separate degree
at most four in a finite product of spins.
\end{lemma}

% Owner: C
\begin{definition}\label{d:pinmodel}
For finite \(G\subset\mathbb Z^2\), let
\[
 \mathcal L(G)=\{(x,z):x\in G,\ z\notin G,\ d(x,z)=1\},\qquad
 S(G)=\{x\in G:d(x,\mathbb Z^2\setminus G)\le2\}.
\]
A pin assignment is \(\tau\in(\mathbb S^2)^{\mathcal L(G)}\).
For \(Y\subset G\), write \(D_Y=G\setminus Y\) and define
\[
 Z^\tau_{D_Y}[f]=
 \sum_{n\in D^{D_Y}}12^{-|D_Y|}f(n)
 \prod_{\substack{\{x,y\}\subset D_Y\\d(x,y)=1}}(1+n_x\cdot n_y)
 \prod_{\substack{(x,z)\in\mathcal L(G)\\x\in D_Y}}
       (1+\tau_{xz}\cdot n_x).
\]
The free sum \(Z^0_{D_Y}[f]\) omits all leg factors.
Write \(Z^\bullet_{D_Y}=Z^\bullet_{D_Y}[1]\),
\(\mathbb E^\bullet_{D_Y}f=Z^\bullet_{D_Y}[f]/Z^\bullet_{D_Y}\),
and \(\Psi_{D_Y}(\tau)=Z^\tau_{D_Y}/Z^0_{D_Y}\).
Here and below \(\bullet\) is either \(\tau\) or \(0\).
The empty partition sum is one. There are no new factors on legs into \(Y\).
Different original legs may carry unrelated unit pins, even if their
outside endpoints coincide.
\end{definition}

% Owner: C
\begin{lemma}\label{l:finitefacts}
Consider a finite network with site laws \(\lambda_{\mathrm U}\) or
\(\lambda_{\mathrm D}\), retained distinct-endpoint bonds
\(1+n_x\cdot n_y\), and retained original unit-pin factors
\(1+\tau_\ell\cdot n_x\), with at most four factors incident to each site.
Its pinned and free partition sums are strictly positive, including
the empty network. Non-neighbours are conditionally independent given
the other spins on every conditioning event of positive probability.
If \(f\) is a product of scalar products of spins and the product of
\(f\) with the bond factors has separate degree at most four at every
site, then \(Z^0[f]\ge0\), hence \(\mathbb E^0 f\ge0\).
All normalized spins have norm at most one.
\end{lemma}

% Owner: C
\begin{lemma}\label{l:35}
For \(0\le k\le3\) and \(|a_j|\le1\), the probability law proportional to
\(\prod_{j=1}^k(1+n\cdot a_j)\lambda_{\mathrm U}(\mathrm dn)\)
has mean of norm at most \(3/5\). The same bound holds with
\(\lambda_{\mathrm D}\) in place of \(\lambda_{\mathrm U}\).
Moreover,
\[
 \left\{
 \frac{\mathcal M\!\left(n\prod_{j=1}^3(1+n\cdot a_j)\right)}
      {\mathcal M\!\left(\prod_{j=1}^3(1+n\cdot a_j)\right)}
       :a_1,a_2,a_3\in\mathbb S^2
 \right\}=\{b\in\mathbb R^3:|b|\le3/5\}.
\]
Consequently an unpinned site having at most three present factors
has conditional mean of norm at most \(3/5\), on every positive
conditioning fibre.
\end{lemma}

% Owner: C
\begin{lemma}\label{l:exact}
At each site of a finite set choose any positive weighted unit-vector
law whose moments through degree four equal those of \(\mathcal M\).
A polynomial of separate degree at most four has the same product
functional under these laws as under \(\lambda_{\mathrm U}\) at all
sites. The analogous assertion holds after replacing any chosen site
laws by \(3/5\) times that law plus \(2/5\,\delta_0\).
In particular this applies to partition products of factors
\(1\pm n_x\cdot n_y\) and \(1\pm n_x\cdot a\), \(|a|=1\), whenever
each variable occurs at most four times, and to inserted observables
only when their degrees are included in this same bound.
\end{lemma}

% Owner: C
\begin{definition}[(P)]\label{h:P}
For \(A\ge2\) and \(0<\theta<1\), property \(\mathsf P(A,\theta)\)
means that, for every finite \(G\subset\mathbb Z^2\), every unit-pin
assignment on \(\mathcal L(G)\), every
\(Y\subset G\setminus S(G)\), and all \(u,v\in G\setminus Y\) with
\(d(u,v)=2\) and a common neighbour in \(Y\),
\[
 \left|\mathbb E^\tau_{G\setminus Y}(n_u\cdot n_v)
             -\mathbb E^0_{G\setminus Y}(n_u\cdot n_v)\right|
 \le A\theta^{\min\{d(u,S(G)),d(v,S(G))\}} .
\]
The distances are ambient plane distances to the original \(S(G)\);
removed sites are not added to the boundary.
\end{definition}


\section{The quantum model, uniqueness, and positive frames}

% Owner: A
\begin{definition}\label{d:poly}
Let \(V_d\) be the homogeneous binary polynomials of degree \(d\).
Its monomials \(z_1^{d-j}z_2^j\), \(0\le j\le d\), are orthogonal with
squared norms \(\binom dj^{-1}\). The action is
\((U_gp)(z)=p(g^{-1}z)\). Use \(V_4\) at each physical site.
For each of the finitely many frame letters \(a\) choose one unit
spinor \(\zeta(a)\in\mathbb C^2\) with
\(\zeta(a)\zeta(a)^*=(I+a\cdot\sigma)/2\).
No spinor choice on the whole sphere is part of the definitions.
\end{definition}

% Owner: A
\begin{definition}\label{d:frame}
A frame is a finite set \(A\subset\mathbb S^2\) with positive weights
\(w(a)\) summing to one such that
\(\sum_{a\in A}w(a)p(a)=\mathcal M(p)\)
for every real polynomial \(p\) of degree at most four.
\end{definition}

% Owner: A
\begin{lemma}\label{l:tight}
For every frame \((A,w)\) and every \(p\in V_4\),
\[
 \sum_{a\in A}5w(a)|p(\zeta(a))|^2=\|p\|_{V_4}^2.
\]
For site frames \((D_x,w_x)\), set \(\Omega=\prod_xD_x\),
\(\pi(s)=\prod_x5w_x(s_x)\), and
\(\iota\varphi(s)=\varphi((\zeta(s_x))_x)\).
Then \(\iota:\mathcal H_\Lambda\to\ell^2(\Omega,\pi)\) is an isometry.
\end{lemma}

% Owner: A
\begin{lemma}\label{l:threeframes}
Let
\[
 T=\begin{pmatrix}99/101&-20/101&0\\20/101&99/101&0\\0&0&1\end{pmatrix},
 \quad D^{(j)}=T^jD\ (j=0,1,2),\quad
 b=\frac{5-\sqrt5}{505},\quad r_b=\frac b{2-b}.
\]
The matrix \(T\) is in \(\mathrm{SO}(3)\).
Each \(D^{(j)}\), with uniform weights \(1/12\), is an antipodal frame.
For \(i\ne j\), \(a\in D^{(i)}\), \(a'\in D^{(j)}\),
\[
 0<b<1,\qquad b\le1\pm a\cdot a'\le2-b,\qquad0<r_b<1.
\]
The common lower bound \(b\) is attained, and distinct frames are
disjoint even up to antipodal signs.
\end{lemma}

% Owner: A
\begin{definition}\label{d:framemodel}
Given any proper three-colouring \(\kappa\) of a plane or torus,
put \(D_x=D^{(\kappa(x))}\), with weights \(1/12\).
On the torus define
\[
 \mu^{\mathrm{af}}_{\kappa,L,N}(s)=
 \frac1{Z^{\mathrm{af}}}\prod_x\frac1{12}
       \prod_{\{x,y\}\in E_{L,N}}(1-s_x\cdot s_y).
\]
On the plane, for finite \(W\subset\mathbb Z^2\) and
\(\eta_z\in D_z\) at exterior neighbour sites, define
\[
 \mu^{+,\eta}_{\kappa,W}(n)=\frac1{Z^{+,\eta}_{\kappa,W}}
 \prod_{x\in W}\frac1{12}
 \prod_{\substack{\{x,y\}\subset W\\d(x,y)=1}}(1+n_x\cdot n_y)
 \prod_{\substack{x\in W,z\notin W\\d(x,z)=1}}(1+n_x\cdot\eta_z).
\]
Only the finitely many needed exterior letters are specified.
These laws have full support. Write \(E_U\) for conditioning off \(U\),
\(\operatorname{Var}_U f=E_U|f-E_Uf|^2\), and
\[
 \mathcal D_\mu(f)=\sum_x\mu(\operatorname{Var}_x f),\qquad
 \lambda(\mu)=\inf_{\operatorname{Var}_\mu f>0}
             \frac{\mathcal D_\mu(f)}{\operatorname{Var}_\mu f}.
\]
All generators have one update per physical site. A gap lower bound is
always equivalently a Poincare inequality, including one-point spaces.
\end{definition}

% Owner: A
\begin{lemma}\label{l:planeenergy}
All the frame specifications above are consistent under conditioning
and have nearest-neighbour Markov dependence. Uniformly over proper
colourings, finite volumes, admissible frame-valued boundaries, and
fully specified one-site fibres, every letter has conditional
probability at least
\[
 \alpha=\frac1{12}r_b^4\in(0,1].
\]
The same lower bound holds after conditioning on fewer letters.
Any prescribed assignment on \(k\) sites, under arbitrary consistent
partial conditioning off those sites, has conditional probability
at least \(\alpha^k\).
For a finite plane block \(Q\), the single-site heat bath of every conditional block law in \(Q\) has Poincare gap at least
\(r_b^{|E(Q)|}\), where \(E(Q)\) is its internal edge set.
\end{lemma}


\section{Windows, stopped coupling, and the periodic block gap}

% Owner: A (frozen shared definition)
\begin{definition}\label{d:shells}
For \(a\in\mathbb Z^2\), integer \(u\ge1\), put
\(C_u(a)=a+\{-u,\ldots,u\}^2\).
For \(\operatorname{dist}_\infty(y,C_u(a))=1\), let
\(s=u^{1/4}\), \(k_i=\lfloor is\rfloor\), \(i_1=\lfloor s\rfloor+1\).
For an integer \(l_0\ge1\) define
\[
 \begin{split}
 F_i&=\{x\in C_u(a):|x-y|_\infty>(i-1)s\},\\
 D_i&=\{x\in C_u(a):|x-y|_\infty=k_i\},\\
 A_i&=\{x\in D_i:\operatorname{dist}_\infty
                    (x,\mathbb Z^2\setminus C_u(a))\le l_0\},\qquad
 B_i=D_i\setminus A_i .
 \end{split}
\]
Let \(I_0\) consist of the integers \(2\le i\le i_1\) for which
some point of \(D_i\) is at distance strictly less than \(s/3\)
from a corner site of \(C_u(a)\).
All these are sets of physical sites.
\end{definition}

% Owner: A (contract); B (establishes it)
\begin{definition}[(W)]\label{d:W}
The family of plane frame specifications has \(\mathsf W(l_0,p_0,M_\star)\)
if \(l_0,M_\star\) are positive integers, \(0<p_0\le1\), and,
for every proper plane three-colouring, every translate \(C_u(a)\)
with integer \(u\ge M_\star\), every \(y\) as in Definition~\ref{d:shells},
every integer \(2\le i\le u^{1/4}\) outside \(I_0\), and every two
admissible frame-valued boundary conditions \(\eta,\xi\) for \(F_i\),
\[
 d_{\mathrm{TV}}\!\left(
       \mu^{+,\eta}_{\kappa,F_i,B_i},
       \mu^{+,\xi}_{\kappa,F_i,B_i}\right)\le1-p_0 .
\]
The subscript \(B_i\) denotes a marginal. Boundaries are arbitrary
within their site alphabets; they need not agree anywhere.
\end{definition}


\section{The cut and the site-scale window theorem}

% Owner: B
\begin{lemma}\label{l:shellgeom}
For \(u\ge M_{\mathrm g}(l_0)=\max\{16,6(l_0+1)\}^4\),
\(|I_0|\le1\), and \(D_i\subset F_i\) for \(2\le i\le i_1\).
For \(2\le i\le i_1\) outside \(I_0\), \(|A_i|\le8l_0\).
Every side of \(y+[-k_i,k_i]^2\) parallel to an edge of \(C_u(a)\)
and meeting \(C_u(a)\) is at distance at least \(s/3\) from the
exterior half-plane beyond that edge.
For \(i\ge2\), every point of \(C_u(a)\setminus F_i\)
within \(\ell^\infty\) distance one of \(F_i\) lies in \(D_{i-1}\);
\(y\) is not within that distance of \(F_i\).
\end{lemma}

% Owner: B
\begin{theorem}\label{t:cutMOS}
If \(\mathsf P(A,\theta)\) holds, then the family of plane frame
specifications has \(\mathsf W(64,p_0,M_\star)\), with common
\(p_0\in(0,1]\) and integer \(M_\star\ge1\) depending only on
\(A,\theta\). The constants work for every proper plane
three-colouring and every translate.
Specifically, let
\[
 E=1+4A+10A\theta^{-4}
       \left[48\theta^{14}\left(17+\frac1{1-\theta}\right)+1\right],
 \qquad p_0=e^{-2E}.
\]
There exists an integer \(M_\star\ge1\) for which (W) holds
with this displayed choice of \(p_0\).
\end{theorem}


\section{Walk expansion without scalar trace masses}

% Owner: C
\begin{definition}\label{d:domain}
Fix finite \(G\), \(S=S(G)\), and original unit pins \(\tau\).
A domain assigns statuses \(\mathrm U,\mathrm D,\mathrm X\);
\(\mathrm X\) sites are removed. Each nonremoved site has a kept
subset of its incident bonds inside \(G\). A bond is present if both
ends are nonremoved and both keep it. A \(\mathrm D\) site has at
most two present bonds. Protected sites in \(S\) have status
\(\mathrm U\) and keep every incident bond inside \(G\).
Use \(\lambda_{\mathrm U},\lambda_{\mathrm D}\) at the corresponding
sites and define
\[
 Z_D^\tau[f]=\sum_n f(n)
       \prod_{\text{present }\{x,y\}}(1+n_x\cdot n_y)
       \prod_{\substack{(x,z)\in\mathcal L(G)\\x\in S}}
                   (1+\tau_{xz}\cdot n_x)
       \prod_{x:\,\mathrm{st}(x)\ne\mathrm X}\lambda_{\mathrm{st}(x)}(n_x).
\]
The sum is over the product of the supports \(D\) at \(\mathrm U\)
sites and \(D\cup\{0\}\) at \(\mathrm D\) sites.
Define the free sum, normalized expectations and \(\Psi_D\) as before.
The class \(\mathcal D_{\mathrm{cons}}(G)\) consists of domains in
which every \(\mathrm U\) site keeps all its bonds, every \(\mathrm D\)
site keeps exactly two, and every bond between nonremoved sites is
kept by both ends or by neither.
An auxiliary network has the same data but permits protected
\(\mathrm U\) sites to have arbitrary kept subsets. It retains every
original leg factor. Neither domains nor auxiliary networks carry
a scalar trace mass. Only the networks occurring below, with at most
four factors at each site, are used.
\end{definition}

% Owner: C
\begin{definition}\label{d:pattern}
Let \(\mathsf T=\{z:d(z,0)\le2\}\), with passage site \(0\),
arrival site \(\mathsf s=(0,-1)\), other neighbours
\(\mathsf e=(1,0),\mathsf N=(0,1),\mathsf W=(-1,0)\),
four corners \((\pm1,\pm1)\), and tips \((\pm2,0),(0,\pm2)\).
The external slots are the two external bonds at each corner and
the three at each tip, twenty in total.
A pattern chooses an out-neighbour \(o\in\{\mathsf e,\mathsf N,\mathsf W\}\)
and statuses/kept sets with the following rules.
The passage site is \(\mathrm U\) and keeps precisely its three
non-arrival bonds. The arrival site is removed or is \(\mathrm D\)
keeping two of its three non-passage bonds.
Each other neighbour is \(\mathrm U,\mathrm X\), or \(\mathrm D\)
keeping its passage bond and one other bond; \(o\) is not removed.
Each corner is \(\mathrm U,\mathrm X\), or \(\mathrm D\) keeping
two of its four bonds. Every \(\mathrm U\) site other than the
passage site keeps all four bonds.
Between nonremoved non-tip sites, a bond is kept by both ends or
by neither. Tips are covering \(\mathrm U\) sites; their bond to
a neighbour \(z\) is present exactly when \(z\) is nonremoved and
keeps that bond. No consistency condition with a global tip kept
set is imposed.
For \(m\in(\mathbb S^2)^{20}\), include every external slot factor
at every nonremoved corner and tip, including both corner slots
when the corresponding actual external bonds are missing.
Let \(Z_{\rm cur}(\pi;m)\) be the finite product functional of
all present template bonds and these slots, using the status laws.
Let \(Z_{\rm next}(\pi;m)\) change the passage law to
\(\lambda_{\mathrm D}\) and drop its out-bond.
The resulting class is \(\mathcal P_{\rm ext}\); it is straight
when \(o=\mathsf N\), and turning otherwise.
\end{definition}

% Owner: C
\begin{definition}[(H1)]\label{h:H1}
For \(c_{\mathrm s},c_{\mathrm t}\in(0,1]\), (H1) means that
for every \(\pi\in\mathcal P_{\rm ext}\) and every
\(m\in(\mathbb S^2)^{20}\),
\[
 Z_{\rm cur}(\pi;m)\ge c_{\rm type(\pi)}Z_{\rm next}(\pi;m).
\]
In the paper the fixed thresholds are
\(c_{\mathrm s}=961/1000\), \(c_{\mathrm t}=187/200\).
No minimum, minimizing pins, or optimal inversion is asserted.
\end{definition}

% Owner: C
\begin{definition}\label{d:weights}
Put \(x_{\mathrm s}=1/c_{\mathrm s}\), \(x_{\mathrm t}=1/c_{\mathrm t}\).
A lattice trail \(w=(v_0,\ldots,v_n)\) has distinct unoriented
nearest-neighbour bonds. Give it a designated unit incoming direction
\(a\), and forbid \(v_1-v_0=-a\) when \(n\ge1\); the virtual
incoming bond is not itself a used bond of the trail.
For departure \(i\), \(0\le i<n\), give weight \(1/5\) if \(v_i\)
has been departed from earlier, and otherwise weight \(x_{\mathrm s}/3\)
for continuing in the incoming direction and \(x_{\mathrm t}/3\)
for a turn. Let \(\mathsf W(w)\) be the product, with
\(\mathsf W(\varnothing)=1\).
\end{definition}

% Owner: C
\begin{definition}[(H2)]\label{h:H2}
For \(C_V\ge1\) and \(0<q<1\), (H2) means that for every
integer \(n\ge0\), every start and every designated incoming direction,
\[
 \sum_{\substack{|w|=n\\\text{initial reversal forbidden}}}
                      \mathsf W(w)\le C_Vq^n .
\]
The sum is over plane lattice trails with those specified initial data.
\end{definition}


\section{Pair closure, a positive-vector witness, and renewal}

% Owner: C
\begin{definition}\label{d:kernel}
For \(X=B,C\), fix \(c=0\), \(u=(-1,0)\), and \(v=(1,0)\)
for \(B\), \(v=(0,1)\) for \(C\).
Let \(\mathcal T_n^X\) be the length-\(n\) plane trails avoiding
\(c\), from \(u\) to their first hit of \(v\).
The initial incoming direction is \(u-c\).
Departure \(i\), \(0\le i<n\), is a lattice revisit if \(v_i\)
was departed from earlier. At each lattice first visit put
\(x_i=x_{\mathrm s}\) or \(x_{\mathrm t}\) according to its type.
Of the three neighbours other than the arrival site, the pair not
containing the outgoing neighbour has coefficient
\(\bar\alpha=18094/100000\), and each other pair has coefficient
\(\bar\beta=36751/100000\). A pair is dead if it contains \(c\) or a
site departed from twice before departure \(i\).
A live pair \(p\) has type \(Y=B,C\) and shift
\[
 s(p)=\max_{z\in p}\min\{d(z,u),d(z,v)\}.
\]
Let \(\kappa_{Y,s}(i)\) sum its coefficients at that type and shift.
Choose once and for all, for each trail \(w\), an ordering
\(\sigma_w\) of its lattice first visits, independent of any later
\(\theta\). Define
\[
 M_{XY,s}=\sum_{n\ge1}\sum_{w\in\mathcal T_n^X}
 3^{-n}(3/5)^{\operatorname{rev}(w)}
 \sum_k\kappa_{Y,s}(\sigma_w(k))\prod_{l\le k}x_{\sigma_w(l)},\qquad
 M_\theta[X,Y]=\sum_{s\ge0}M_{XY,s}\theta^{-s}.
\]
All sums are nonnegative series. The identity
\[
 3^{-n}(3/5)^{r}=3^{-(n-r)}5^{-r}
\]
identifies these coefficients with the new trace convention when
the statuses are determined by the lattice trail.
\end{definition}

% Owner: C
\begin{definition}[(H3\( _h\))]\label{h:H3}
For the same \(c_{\mathrm s},c_{\mathrm t}\) and the kernel above,
(H3\( _h\)) means that there exist fixed orderings \(\sigma_w\),
a vector \(h=(h_B,h_C)\in\mathbb Q_{>0}^2\), and
\(\rho_0\in\mathbb Q\cap(0,1)\) such that, for \(X=B,C\) and
every integer \(R\ge0\),
\[
 \sum_{n=1}^{R}\sum_{w\in\mathcal T_n^X}
 3^{-n}(3/5)^{\operatorname{rev}(w)}
 \sum_k\left(\sum_{Y,s}\kappa_{Y,s}(\sigma_w(k))h_Y\right)
                          \prod_{l\le k}x_{\sigma_w(l)}
 \le\rho_0h_X .
\]
Equivalently the full nonnegative sums obey \(M_1h\le\rho_0h\).
An ordering is only a supplied witness; no optimality is required.
\end{definition}

% Owner: C
\begin{lemma}\label{l:Mtheta}
Assume (H2) and (H3\( _h\)).
For every \(q^2<\theta\le1\), all entries of \(M_\theta\) are
finite. There exist \(q^2<\theta<1\) and \(\rho_0<\rho<1\)
such that the same vector \(h\) satisfies \(M_\theta h\le\rho h\).
Every shift of a length-\(n\) trail is at most \(n/2+1\), and
each matrix entry is bounded by
\[
 \sum_{n\ge1} n(\bar\alpha+2\bar\beta)
                      C_V q^n\theta^{-n/2-1}.
\]
\end{lemma}

% Owner: C
\begin{theorem}\label{t:rec}
Assume (H1), (H2), and (H3\( _h\)) at the same pair
\((c_{\mathrm s},c_{\mathrm t})\).
Then there exist \(A\ge2\) and \(0<\theta<1\) such that
\(\mathsf P(A,\theta)\) holds.
One may take the \(\theta,\rho,h\) supplied by the shifted-witness
lemma and
\[
 A=\frac{\max h}{\min h}\,
       \frac{\max\{2,\,2(12/5)^{24}(C_V/(1-q))^2\}}{1-\rho}.
\]
\end{theorem}


\section{Exact statements of the computational inputs}

% Owner: C
\begin{theorem}\label{t:H1}
For every \(\pi\in\mathcal P_{\rm ext}\) and every
\(m\in(\mathbb S^2)^{20}\),
\[
 Z_{\rm cur}(\pi;m)\ge
 \begin{cases}
 (961/1000)Z_{\rm next}(\pi;m),&\pi\text{ straight},\\
 (187/200)Z_{\rm next}(\pi;m),&\pi\text{ turning}.
 \end{cases}
\]
\end{theorem}

% Owner: C (theorem text and proof account); separate certificate owner supplies H2/H3 witnesses
\begin{theorem}\label{t:inputs}
At \(c_{\mathrm s}=961/1000\) and \(c_{\mathrm t}=187/200\),
(H1) holds. There exist rational constants
\(C_V\ge1\), \(0<q<1\), \(h_B,h_C>0\), \(0<\rho_0<1\),
and fixed orderings \(\sigma_w\), for which (H2) and
(H3\( _h\)) hold with these constants and orderings.
\end{theorem}


\end{document}
